% vim: ai sw=2 spell spelllang=cs

\begin{frame}{UNIX a síť}

  \begin{itemize}
    \item Na UNIXu býval všechno soubor

    \pause

    \item Na soubor existuje souborový deskriptor

    \pause

    \item I na síťové spojení se vytváří souborový deskriptor
    \begin{itemize}
      \item Tzv.\ socket (\doc{man 7 socket})
    \end{itemize}

    \pause

    \item Na něm se dá provádět většina standardních operací
  \end{itemize}

\end{frame}

\begin{frame}[fragile]{\code{socket}}

  \begin{itemize}
    \item Otevření: \code{int socket(int domain, int type, int protocol)}
    \begin{itemize}
      \item Hlavičky: \header{sys/types.h}, \header{sys/socket.h}
      \item \code{domain}: protokol -- \code{AF_UNIX} \code{AF_INET},
	\code{AF_INET6}, \dots
      \item \code{type}: \code{SOCK_STREAM} (TCP), \code{SOCK_DGRAM} (UDP), \dots
      \item \code{protocol}: pokud je na výběr (NETLINK), jinak 0 (pro UDP i TCP)
      \item Návratová hodnota: souborový deskriptor nebo \code{\-1} při chybě
    \end{itemize}

    \pause

    \item Zavření
    \begin{itemize}
      \item Klasicky: \code{close}
\ifpcdirlinuxb
      \item Částečně: \code{int shutdown(int sockfd, int how)}
      \begin{itemize}
	\item \code{how} -- \code{SHUT_RD}, \code{SHUT_WR}, \code{SHUT_RDWR}
      \end{itemize}
\fi
    \end{itemize}

    \pause
  \end{itemize}

  \onslide<1->

  \begin{block}{Příklad (bez ošetření chyb)}
  \begin{lstlisting}
int s = socket(AF_UNIX, SOCK_STREAM, 0);
close(s);
  \end{lstlisting}
  \end{block}

\end{frame}
